\section{Consequences for the rationally connected quotient}
\label{sec:consequences}

The main theorem also applies to the general fiber of the rationally
connected quotient of a uniruled compact K\"ahler manifold. That fiber
is a smooth projective manifold even when the ambient manifold is not
projective. Its tangent splitting is therefore integrable, which supplies
the premise in H\"oring's structure theorems \cite{Hor08}.

\begin{corollary}\label[corollary]{cor:rational-quotient}
Let \(X\) be a connected uniruled compact K\"ahler manifold with a
holomorphic splitting \(T_X=V_1\oplus V_2\). Let \(Z\) be a general
fiber of its rationally connected quotient and set
\(W_i=T_Z\cap V_i|_Z\).
For each \(i=1,2\), there is an almost-holomorphic fibration
\(\phi_i:X\dashrightarrow Y_i\) whose general fiber \(F_i\) is rationally
connected and, for compatible choices \(F_i\subset Z\), satisfies
\[
T_{F_i}=W_i|_{F_i}\subset V_i|_{F_i}.
\]
If \(W_i=0\), one may take \(\phi_i=\id_X\), with point fibers.

If \(\rk V_1=2\), exactly one of the following holds:
\begin{enumerate}[label=(\roman*)]
\item \(W_1=V_1|_Z\), and \(X\) admits an analytic fiber bundle
\(X\to Y\) with rationally connected general fiber and \(T_{X/Y}=V_1\).
\item \(W_1\) is a line bundle, and there is an equidimensional
holomorphic map \(X\to Y\) whose general fiber \(F\) is a rational curve
with \(T_F\subset V_1|_F\).
\item \(W_1=0\), equivalently \(T_Z\subset V_2|_Z\).
\end{enumerate}
\end{corollary}

\begin{proof}
Use the published version of \cite{Hor08}; its definition of an
almost-holomorphic fibration includes the identity case. Choose \(Z\)
over a regular value in the fibration locus of the quotient supplied by
\cite[Definition~2.9 and Theorem~2.10]{Hor08}. Then \(Z\) is an actual
smooth connected compact submanifold of \(X\). It inherits a K\"ahler
form and is rationally connected. Campana's theorem, recalled in the
remark after \cite[Definition~2.8]{Hor08}, makes \(Z\) projective.

By \cite[Proposition~3.16]{Hor08}, the ambient complementary projections
preserve \(T_Z\). Their restrictions are complementary holomorphic
idempotents. An idempotent has locally constant rank, so on connected
\(Z\) their images are the holomorphic subbundles \(W_i\), and
\[
T_Z=W_1\oplus W_2.
\]
If both ranks are positive, then \(\dim Z\ge2\), and \cref{thm:main}
makes both \(W_i\) integrable. Otherwise they are \(0\) and \(T_Z\),
which are integrable as well.

This is the general-fiber integrability hypothesis of
\cite[Theorem~1.3]{Hor08}, which gives the fibrations \(\phi_i\) and the
stated tangent identities. Its proof includes the identity when
\(W_i=0\). If \(\rk V_1=2\), \cite[Theorem~1.4]{Hor08} gives the three
conclusions according as \(\rk W_1\) is \(2\), \(1\), or \(0\).
These ranks are mutually exclusive and exhaustive, proving the stated
trichotomy.
\end{proof}
